\documentclass[aspectratio=169,10pt]{beamer}
\usepackage[UTF8,fontset=none]{ctex}
\usepackage{booktabs}
\usepackage{array}
\usetheme{tutorial-showcase}

\title[VSP-Beamer Showcase]{VSP-Beamer Showcase}
\subtitle{渐变封面 / 章节导航 / 进度页脚 / 数据组件}
\VSPsetspeaker{VSP-Beamer Team}{VSP-Beamer}
\date{2026年11月}
\VSPsetupfootline{VSP-Beamer Showcase}{2026}{VSP-Beamer Team}

\begin{document}
\VSPtitleframe

\section{页面骨架}

\begin{frame}[t]{标题区与页脚}
\begin{columns}[T,onlytextwidth]
  \begin{column}{.56\textwidth}
    标题上方的小字是\VSPmark{章节面包屑}，由 \texttt{\textbackslash section} 自动生成；标题下方是红到青的渐变分隔线。

    页脚最底部是\VSPmark{全局进度条}，中间的圆点对应全部章节：已讲完的为红色，当前章节为青色。

    \begin{itemize}
      \item 章节页自带议程：右侧列出全部章节并高亮当前位置
      \item 封面、章节页与尾页共用深红渐变与丝带曲线
      \item \alert{alert} 为主题红加粗，\textbf{textbf} 为标准加粗
    \end{itemize}
  \end{column}
  \begin{column}{.40\textwidth}
    \begin{vspquote}{使用方法}
      \texttt{\textbackslash usetheme\{tutorial-showcase\}}

      其余接口与 \code{tutorial-red} 相同，文稿可以直接切换。
    \end{vspquote}
    \VSPchip{tutorial-red} \VSPchip{report-red} \VSPchip[VSPTeal]{tutorial-showcase}
  \end{column}
\end{columns}
\end{frame}

\section{数据与流程}

\begin{frame}[t]{关键数字}
\begin{columns}[T,onlytextwidth]
  \begin{column}{.235\textwidth}\VSPstat{7}{主题入口}{共享一套实现}\end{column}
  \begin{column}{.235\textwidth}\VSPstat[VSPTeal]{16:9}{页面比例}{aspectratio=169}\end{column}
  \begin{column}{.235\textwidth}\VSPstat[VSPCalloutBlue]{XeLaTeX}{编译引擎}{latexmk 驱动}\end{column}
  \begin{column}{.235\textwidth}\VSPstat[VSPGreen]{100\%}{离线构建}{无远程图片}\end{column}
\end{columns}

\vspace{4pt}
\texttt{\textbackslash VSPstat[颜色]\{数字\}\{标签\}\{备注\}} 适合放在多栏中并排展示关键数字。

\VSPmeter{0.82}{进度条 \texttt{\textbackslash VSPmeter}}{82\%}
\VSPmeter[VSPTeal]{0.46}{可指定颜色}{46\%}
\end{frame}

\begin{frame}[t]{流程条与提示框}
\VSPsteps{提出问题,建立方法,实现系统,实验评估,总结展望}

\vspace{6pt}
\begin{columns}[T,onlytextwidth]
  \begin{column}{.32\textwidth}
    \begin{vspquote}{强调引用}
      彩色标题压在一条细线上，正文不带底色和边框。
    \end{vspquote}
    \begin{vspquoteyellow}{黄色提示}
      适合补充说明。
    \end{vspquoteyellow}
  \end{column}
  \begin{column}{.32\textwidth}
    \begin{vspquoteblue}{蓝色提示}
      适合方法或定义。
    \end{vspquoteblue}
    \begin{vspquotegreen}{绿色提示}
      适合已完成的结果。
    \end{vspquotegreen}
  \end{column}
  \begin{column}{.32\textwidth}
    \begin{vspquotepurple}{紫色提示}
      适合旁注与延伸。
    \end{vspquotepurple}
    \begin{block}{Beamer block}
      标准 \code{block} 同样适用。
    \end{block}
  \end{column}
\end{columns}
\end{frame}

\section{图表与代码}

\begin{frame}[t]{图文与表格}
\begin{columns}[T,onlytextwidth]
  \begin{column}{.36\textwidth}
    \begin{vspquote}{图片}
      \begin{Figure}[奶龙：\code{nailong-wave.png}]
        \includegraphics[height=3.2cm]{\VSPasset{nailong-wave.png}}
      \end{Figure}
    \end{vspquote}
  \end{column}
  \begin{column}{.60\textwidth}
    \begin{Table}[主题对照]
      \footnotesize
      \begin{tabular}{@{}lccc@{}}
        \toprule
        主题 & 配色 & 封面 & 页脚 \\
        \midrule
        \texttt{tutorial-red} & 红 & 白底标题 & 三格信息栏 \\
        \texttt{tutorial-red-shtu} & 红 & 校园背景 & 页码 \\
        \texttt{report-red} & 红 & 居中标题块 & 三格信息栏 \\
        \texttt{tutorial-showcase} & 红 + 青 & 渐变与丝带 & 进度条 \\
        \bottomrule
      \end{tabular}
    \end{Table}
  \end{column}
\end{columns}
\end{frame}

\begin{frame}[t,fragile]{代码与公式}
\begin{columns}[T,onlytextwidth]
  \begin{column}{.48\textwidth}
    \begin{enumerate}
      \item 列出全部文稿
      \item 严格检查全部文稿
    \end{enumerate}
\begin{lstlisting}[language=bash]
make list
make check
\end{lstlisting}
  \end{column}
  \begin{column}{.48\textwidth}
    \begin{vspquoteblue}{行间公式}
      \[
        \mathcal{L} = \sum_{i=1}^{N} \bigl\lVert \hat{I}_i - I_i \bigr\rVert_1
          + \lambda\,\mathcal{L}_{\mathrm{SSIM}}
      \]
    \end{vspquoteblue}
  \end{column}
\end{columns}
\end{frame}

\VSPendframe{VSP-Beamer \textcolor{white!60}{/} Showcase}
\end{document}
